\section{Soundness of the certificate interface}
\label{sec:soundness}
This section proves the mathematical implications used by the checker. The results apply to accepted evidence satisfying their hypotheses; they do not assert that every convex MINLP has evidence accepted by the implementation. We separate the nonlinear support calculation, the discrete proof, and their composition so that a defect in one interface cannot be hidden by a successful check of another.

\subsection{Rational propagation preserves feasible points}
\label{sec:propagation}
Start with the declared box, intersect binary coordinates with $[0,1]$, and substitute exact fixed values. For a row $\sum_j c_jx_j\leq u$, let $R_k^-$ be a lower bound on $\sum_{j\ne k}c_jx_j$ over the current box. If $c_k\ne0$ and $R_k^-$ is finite, the implied bound is
\begin{equation}
 x_k\leq (u-R_k^-)/c_k\quad(c_k>0),\qquad
 x_k\geq (u-R_k^-)/c_k\quad(c_k<0).
 \label{eq:propagate-upper}
\end{equation}
For a lower row $\sum_jc_jx_j\geq\ell$, use an upper bound $R_k^+$ for the other terms and divide $\ell-R_k^+$ by $c_k$, reversing the corresponding directions. The endpoints of each term are obtained by its coefficient sign. A missing endpoint gives no finite tightening in that direction. All finite operations are rational. An upper bound on an integer coordinate may be rounded down; a lower bound may be rounded up.

\begin{lemma}[Feasible-point preservation]
\label{lem:propagation}
Any finite sequence of these operations produces a box $B$ containing every original mixed-integer feasible point. If a coordinate interval becomes empty, the original feasible set is empty.
\end{lemma}
\begin{proof}
Initially every feasible point satisfies the declared, fixed, and binary bounds. Suppose it lies in the box before an update. Its other terms are at least $R_k^-$, so the upper row implies $c_kx_k\leq u-R_k^-$. Division gives \eqref{eq:propagate-upper}. The lower-row argument is identical with an upper bound on the other terms. Integer rounding preserves each feasible integral coordinate. Intersecting with an implied bound therefore preserves the point. Induction proves the inclusion; an empty coordinate interval cannot contain such a point.
\end{proof}
The implementation makes at most 20 sweeps. This is a strength limit, not a soundness assumption. Bounds from integer rounding need not contain the continuous relaxation: the inclusion needed here is $F\subseteq B$. Every cut is nevertheless justified on the whole resulting box. A constant affine row is either a tautology or a contradiction. The finite-bound interface omits the former and rejects the latter; it does not issue a nonlinear infeasibility certificate merely because propagation detects a contradiction.

\subsection{Safe rational affine underestimators}
\label{sec:safe-cuts}
The finite-box calculation below is support minimization applied to an affine residual, an established rigorous-bounding principle \citep{jansson2004-rigorous-convex,messine2019-reliable-relaxation}. We state it with explicit rational-enclosure and unbounded-coordinate conditions because these are the replay interface.
Write a normalized row, including an objective epigraph row if necessary, as
\begin{equation}
 r(x)=g(x)+q^\top x+c.
 \label{eq:row-decomposition}
\end{equation}
Here $q,c$ are rational and $g$ is convex on the certified box. Let $z\in B$ and suppose a finite vector $p$ satisfies
\begin{equation}
 g(x)\geq g(z)+p^\top(x-z)\quad(x\in B).
 \label{eq:g-support}
\end{equation}
A justified subgradient gives this inequality. The implementation uses finite symbolic derivatives in supported differentiable cases; it has no general subgradient oracle. Rationality is required of the supplied point $z$ and cut slope $a$, but not of the true values $g(z)$ or $p$.

Set $\phi(z)=r(z)-a^\top z$ and $d=p+q-a$. For each coordinate define
\begin{equation}
 W_j(d_j)=\sup_{y\in B_j} d_j(z_j-y).
 \label{eq:worst-shift}
\end{equation}
Because $z_j\in B_j$, $W_j\geq0$. The exact formulas, including the conditions for finiteness, are in Table~\ref{tab:shift}. A sign failure in a half-line, or a nonzero residual slope on a free coordinate, makes this particular support bound infinite. This need not mean that the nonlinear function has no finite affine underestimator.

\begin{table}[htbp]
\centering
\caption{Exact coordinate correction, with $z_j\in B_j$. Infinite endpoints are not used in arithmetic.}
\label{tab:shift}
\begin{tabular}{lll}
\toprule
$B_j$ & Condition for finite $W_j$ & $W_j(d_j)$\\
\midrule
$[L_j,U_j]$ & $L_j,U_j$ finite & $\max\{d_j(z_j-L_j),d_j(z_j-U_j)\}$\\
$[L_j,+\infty)$ & $d_j\geq0$ & $d_j(z_j-L_j)$\\
$(-\infty,U_j]$ & $d_j\leq0$ & $d_j(z_j-U_j)$\\
$\R$ & $d_j=0$ & $0$\\
\bottomrule
\end{tabular}
\end{table}

\begin{theorem}[Rational cut safety]
\label{thm:safecut}
Let $B$ be nonempty, $z\in B$, and suppose \eqref{eq:g-support} holds. If every $W_j$ is finite and $b\in\Q$ satisfies
\begin{equation}
 b\leq\phi(z)-\sum_jW_j(d_j),
 \label{eq:safe-intercept}
\end{equation}
then $a^\top x+b\leq r(x)$ for every $x\in B$. In particular, $a^\top x+b\leq0$ is valid for every point in $B$ satisfying $r(x)\leq0$.
\end{theorem}
\begin{proof}
Adding the affine terms to \eqref{eq:g-support} gives
\[
 r(x)-a^\top x\geq \phi(z)+\sum_jd_j(x_j-z_j)
 \geq\phi(z)-\sum_jW_j(d_j)\geq b.
\]
The middle inequality is the definition of each supremum. On a finite interval its value is the larger endpoint value; on a half-line it is finite exactly for the listed sign; on $\R$ only zero slope is bounded above. This also proves Table~\ref{tab:shift}.
\end{proof}

\begin{corollary}[Finite rational enclosure test]
\label{cor:enclosure}
Suppose $v\in\Q$ satisfies $v\leq\phi(z)$ and finite rational intervals $[d_j^-,d_j^+]$ enclose $d_j$. The conditions and quantities in Table~\ref{tab:enclosure} imply $W_j(d_j)\leq\widehat W_j$. Therefore
\begin{equation}
 b\leq v-\sum_j\widehat W_j
 \label{eq:rational-enclosure-test}
\end{equation}
certifies the cut by Theorem~\ref{thm:safecut}.
\end{corollary}
\begin{proof}
On a finite interval, $z_j-L_j\geq0$ and $z_j-U_j\leq0$. The respective products increase with $d_j$ and decrease with $d_j$, giving the two endpoint bounds. The same signs give the half-line bounds, and the sign tests establish finiteness. On a free coordinate the enclosing interval must be exactly $[0,0]$. Finally $v\leq\phi(z)$ and $\widehat W_j\geq W_j$ imply \eqref{eq:safe-intercept}.
\end{proof}

\begin{table}[htbp]
\centering
\caption{Rational sufficient tests for interval gradient data. All intervals must be nonempty and finite.}
\label{tab:enclosure}
\begin{tabular}{lll}
\toprule
$B_j$ & Acceptance condition & $\widehat W_j$\\
\midrule
$[L_j,U_j]$ & $d_j^-\leq d_j^+$ & $\max\{d_j^+(z_j-L_j),d_j^-(z_j-U_j)\}$\\
$[L_j,+\infty)$ & $d_j^-\geq0$ & $d_j^+(z_j-L_j)$\\
$(-\infty,U_j]$ & $d_j^+\leq0$ & $d_j^-(z_j-U_j)$\\
$\R$ & $d_j^-=d_j^+=0$ & $0$\\
\bottomrule
\end{tabular}
\end{table}
For a fixed coordinate, $L_j=U_j=z_j$, so the correction is zero. Exact elimination of a fixed variable avoids differentiating it; without elimination the stated support data must still be justified. Every variable in the row participates, including variables occurring only in $q^\top x$. An absent coefficient in the sparse cut means $a_j=0$, not an absent residual term.

Direct outward interval evaluation of the right-hand side of \eqref{eq:safe-intercept} also gives a valid sufficient test. The implementation uses this variant, with rational input conversion and exact extraction of the final interval's lower endpoint. Dependency loss can weaken the cut but cannot invalidate an actual enclosure. On an upper-unbounded coordinate the producer can choose $a_j$ below a rigorous lower bound for $p_j+q_j$; on a lower-unbounded coordinate it can choose it above a rigorous upper bound. Free coordinates require an exact residual zero. The epigraph coordinate has coefficient $-1$ and zero nonlinear derivative, so taking its slope as $-1$ meets that condition without bounding the epigraph variable.

\begin{example}[Why the intercept must change when the slope is rounded]
\label{ex:rounding}
Let $r(x)=x^2-1/9$ on $[0,1]$ and $z=1/3$. Its tangent is $(2/3)x-2/9$. Replacing the slope by $a=67/100$ while retaining intercept $-2/9$ makes the cut left-hand side equal to $1/900$ at the feasible point $z$. Rounding has excluded an original feasible point. Here $d=2/3-67/100=-1/300$, $\phi(z)=-67/300$, and
\[
 W=\max\{-1/900,1/450\}=1/450.
\]
The safe intercept $b=-203/900$ follows from \eqref{eq:safe-intercept}. Indeed,
\[
 r(x)-\bigl((67/100)x-203/900\bigr)
 =(x-67/200)^2+799/360000>0.
\]
The largest intercept for this fixed slope is $-1/9-(67/100)^2/4$, so the support correction is sufficient without being best possible. If the box were instead $[0,+\infty)$, the slope $67/100$ would fail the residual-sign test. Choosing $a=3/5$ gives $d=1/15$, $W=1/45$, and the valid half-line intercept $b=-2/9$. Thus bounded and unbounded coordinates require different rounding decisions. Neither construction requires $z$ to be an optimizer or an integer point.
\end{example}
Global underestimation is stronger than the implication needed for a feasible-set cut. For example, on $[-1,1]$ the row $x\leq0$ implies $2x\leq0$, although $2x\leq x$ fails for positive $x$. This interface deliberately uses the stronger condition because its evidence is simple to replay.

\subsection{What the discrete proof must establish}
\label{sec:discrete-soundness}
Let $M$ denote the rational mixed-integer linear master, and let $c_M^\top y$ be its linear objective. An objective constant is represented, if needed, by a separate variable fixed to one. We give the argument for minimization. A row is an equality or a weak linear inequality with rational coefficients. A row \emph{dominates} another when it logically implies it by one of the checker's admitted tests: the same left-hand coefficients with a stronger right-hand side and compatible direction, or a constant contradiction. Equalities can dominate either corresponding weak inequality.

A derived row $C$ carries a finite set $A(C)$ of assumption-row indices. Suppose every supplied solution is exactly master feasible, and let $b_M$ be the least objective value of those solutions when any are supplied. Define
\begin{equation}
 S=\begin{cases}
 \{y\in M:c_M^\top y\leq b_M\},&\text{if a solution is supplied},\\
 M,&\text{otherwise}.
 \end{cases}
 \label{eq:incumbent-set}
\end{equation}
A solution cutoff is valid on $S$, not necessarily throughout $M$. The following rules implement this distinction, following the ordinary VIPR solution and assumption semantics \citep{cheung2017-vipr,wood2026-vipr-smt}.

\begin{proposition}[Inference invariant]
\label{prop:vipr-invariant}
Suppose every inference reference points to a strictly earlier row. For every row derived by the rules below, every $y\in S$ satisfying all assumptions in $A(C)$ satisfies $C$.
\begin{enumerate}
\item Original master rows have no assumptions. An assumption row has its own index as its assumption set.
\item A solution row must be dominated by $c_M^\top y\leq b_M$ and requires a supplied feasible solution. Its assumption set is empty relative to $S$.
\item A linear combination uses rational multipliers whose nonzero inequality terms all have the same direction after multiplication; equality terms are unrestricted. Its sum must dominate the claimed row. Its assumptions are the union from all nonzero terms.
\item A rounding step first forms such a combination. Every nonzero left-hand coefficient must be integral and multiply an integer variable. An upper right-hand side is rounded down, or a lower right-hand side up. The rounded row must dominate the claim. Equality rounding is not admitted.
\item An unsplit step references two actual assumption rows $D_1,D_2$ of the form $a^\top y\leq k$ and $a^\top y\geq k+1$, in either order, where $k$ and the linear form are integral on every master-feasible point. It also references branch conclusions $C_1,C_2$ that each dominate $C$. Its assumption set is
\[
 A(C)=\bigl(A(C_1)\setminus\{D_1\}\bigr)
       \cup\bigl(A(C_2)\setminus\{D_2\}\bigr).
\]
Here the notation $\{D_i\}$ means the index of that assumption row.
\end{enumerate}
\end{proposition}
\begin{proof}
Proceed in derivation order. Original rows and assumptions satisfy the invariant immediately. Solution rows hold by \eqref{eq:incumbent-set}. For a combination, a point satisfying the union of dependencies satisfies every nonzero source row. Multiplication with the checked signs and addition give the summed relation, hence its dominated claim. Zero multipliers contribute neither an inequality nor an assumption. For rounding, the left-hand side is an integer at the point, so floor or ceiling of the right-hand side preserves validity.

For an unsplit step, take $y\in S$ satisfying $A(C)$. The integral disjunction is exhaustive, so $y$ satisfies at least one $D_i$. Together with $A(C_i)\setminus\{D_i\}$, this gives every dependency of $C_i$. The induction hypothesis proves $C_i$, which implies $C$. Dependencies involving the other branch or any unrelated assumption remain in the displayed union; discharging one branch cannot erase them from the other branch's proof.
\end{proof}

\begin{theorem}[Unconditional master bound]
\label{thm:master-bound}
Suppose all solutions and every derivation pass the preceding rules, and some assumption-free derived row dominates $c_M^\top y\geq\beta$ for finite $\beta$. Then $c_M^\top y\geq\beta$ for every $y\in M$. No attainment assumption is required. With no supplied solutions, an assumption-free derived contradiction instead proves $M=\varnothing$.
\end{theorem}
\begin{proof}
Proposition~\ref{prop:vipr-invariant} proves the bound on $S$. With no supplied solutions, $S=M$. Otherwise a supplied solution attaining the finite minimum $b_M$ among the finite list belongs to $S$, so $b_M\geq\beta$. Every point of $M\setminus S$ has objective greater than $b_M$, hence at least $\beta$. The contradiction statement follows from the same invariant with $S=M$.
\end{proof}
Maximization uses the greatest checked solution value, a cutoff $c_M^\top y\geq b_M$, and a final upper inequality; the same proof reverses the objective order. An infeasibility claim requires an empty solution section. In particular, a cutoff cannot manufacture infeasibility by excluding feasible solutions. The stronger integer-objective cutoff $b_M-1$ is outside the admitted rules. A valid earlier bound can suffice even when later valid rows are weaker or have unused assumptions; every later row and the end of the file must still be checked. The argument establishes the semantic statement needed by the nonlinear layer, rather than merely matching a numerical endpoint in a proof header.

\subsection{Transfer and primal completion}
\label{sec:bound-transfer}
Let $s=1$ for an original minimization and $s=-1$ for an original maximization, and set $h=s f_{\rm orig}$. Use the exact affine objective if $h$ is affine. Otherwise append a free real coordinate $t$, minimize $t$, and treat $h(x)-t\leq0$ as another convex row. The master includes original affine rows, justified bounds, the original integrality conditions, and accepted cuts. Its extra constant variable, if present, is fixed to one. It has no other constraints without a separate validity justification.

\begin{theorem}[Nonlinear bound transfer]
\label{thm:transfer}
Suppose the box contains $F$, every master cut satisfies Theorem~\ref{thm:safecut} on its applicable box, the proof input is the stated master under a checked variable bijection, and Theorem~\ref{thm:master-bound} establishes its bound $\beta$. Then
\[
 s f_{\rm orig}(x)\geq\beta\quad(x\in F).
\]
Thus $\beta$ is a lower bound for original minimization, and $-\beta$ an upper bound for original maximization. A verified infeasible master would imply $F=\varnothing$.
\end{theorem}
\begin{proof}
For $x\in F$, define $E(x)$ by preserving every original variable, setting $t=h(x)$ if needed, and setting the constant variable to one if present. Original affine rows, integrality, and propagated bounds hold by construction and Lemma~\ref{lem:propagation}. Every nonlinear row is nonpositive at $x$, and every epigraph row is zero at $E(x)$; safe underestimation therefore proves all master cuts there. Thus $E(x)\in M$. Its master objective equals $h(x)$, including the affine constant exactly once. The master bound applied to $E(x)$ proves the result. If $M$ were empty, no such extension of a feasible $x$ could exist.
\end{proof}
For minimization this proves $\beta\leq\inf_{x\in F}f(x)$, including the convention $\inf\varnothing=+\infty$. The proof does not require a minimizer, a Slater point, or a finite optimum before the certificate is supplied. An unbounded relaxation or unsuccessful check provides no finite bound. The mathematical infeasibility implication is broader than the implemented public API, which reports finite bounds only.

\begin{corollary}[Primal completion]
\label{cor:primal}
Suppose an independently checked $\bar x\in F$ has a justified finite objective upper enclosure $f(\bar x)\leq U$ for minimization. A checked lower bound $\beta$ proves
\[
 \beta\leq p^*\leq f(\bar x)\leq U,
 \qquad 0\leq f(\bar x)-p^*\leq U-\beta.
\]
If $U=\beta$, then $\bar x$ attains the exact optimum. The corresponding assertion for maximization follows by sign normalization.
\end{corollary}
\begin{proof}
Feasibility gives $p^*\leq f(\bar x)$ and Theorem~\ref{thm:transfer} gives the other lower bound. The two inequalities give the gap and, in the equality case, attainment.
\end{proof}
The primal check must include variable bounds and integrality as well as rows and objective evaluation. A master solution only supplies the premise used in \eqref{eq:incumbent-set}. A recorded nonlinear reference value supplies neither a feasible point nor a rigorous objective enclosure.
